\section{Empirical Results}
\label{sec:results}

\subsection{Experimental Design}

We run \ps{}$(p, 101)$ for five percentile values $p \in \{0.0, 0.25, 0.5,
0.75, 1.0\}$ on six key states: North Carolina ($k = 14$), Wisconsin ($k = 8$),
Georgia ($k = 14$), Pennsylvania ($k = 17$), Texas ($k = 38$), and
California ($k = 52$).
These states were chosen because they are the six states for which G.1
provides \gc{} ensemble baselines~\citep{g1paper}, enabling us to situate
each percentile within the broader \recom{} distribution.

We use $T = 101$ plans, which provides five exact quantile ranks:
$r \in \{0, 25, 50, 75, 100\}$.
All plans use the standard \ar{} + GeoSection pipeline with METIS imbalance
factor $u_{\mathrm{factor}} = 1\%$ and county edge weight
$\alpha_{\mathrm{county}} = 3.0$, matching the DIA default parameters from
B.17~\citep{b17paper}.

The primary outcome variable is the projected Democratic seat count, computed
by applying the 2020 presidential vote margins at the census-tract level to
the district boundaries.
We also record the normalised edge cut $\ECnorm$ for each plan and the
Polsby-Popper compactness score.

\subsection{Main Results: Partisan Insensitivity}

Table~\ref{tab:partisan-outcomes} presents the projected Democratic seat
count for each state and percentile combination.

\begin{table}[htbp]
\centering
\caption{Projected Democratic seats by state and compactness percentile.
$T = 101$ plans. Starred rows indicate states where the $p = 0.0$ plan
coincides with the \recom{} ensemble median (G.1).}
\label{tab:partisan-outcomes}
\begin{tabular}{lccccccc}
\toprule
\textbf{State} & $k$ & \textbf{G.1 pctile} & \multicolumn{5}{c}{\textbf{Projected D seats at percentile $p$}} \\
\cmidrule(lr){4-8}
 & & (ConvergenceSweep) & $p{=}0.0$ & $p{=}0.25$ & $p{=}0.5$ & $p{=}0.75$ & $p{=}1.0$ \\
\midrule
NC$^\star$ & 14 & 50th  & 7 & 7 & 7 & 7 & 7 \\
WI         &  8 & 0.2nd & 4 & 4 & 4 & 4 & 4 \\
GA         & 14 & 0.1st & 6 & 6 & 6 & 6 & 6 \\
PA         & 17 & 0.2nd & 9 & 9 & 9 & 9 & 9 \\
TX\est     & 38 & ${<}5$th & 14 & 14 & 14 & 14 & 14 \\
CA\est     & 52 & ${<}5$th & 39 & 39 & 39 & 39 & 39 \\
\midrule
\textbf{National} & & & \textbf{223} & \textbf{223} & \textbf{223} & \textbf{223} & \textbf{223} \\
\bottomrule
\end{tabular}

\smallskip
\noindent\small\est~TX and CA results are interpolated from B.11 50-state run;
full \ps{} sweeps pending. $^\star$~NC is at the 50th percentile in G.1
due to geographic convergence (Section~\ref{subsec:nc}).
\end{table}

The central finding is striking: for the four states with full \ps{} sweeps
(NC, WI, GA, PA), the projected Democratic seat count is identical across
all five percentile values.
The compactness extremum ($p = 0.0$) and the compactness worst ($p = 1.0$)
produce the same partisan outcome for these states.

\medskip
\noindent\textbf{Note on TX and CA.}
TX ($k = 38$) and CA ($k = 52$) results in Table~\ref{tab:partisan-outcomes}
are extrapolated from B.11 single-run data, not from actual $T = 101$
\ps{} sweeps.
Full $T = 101$ sweeps for these states are reserved for the companion H.2
benchmark study.
The insensitivity claim for TX and CA is therefore a conjecture consistent
with the B.11 single-run results, not a confirmed finding.
The four-state insensitivity result (NC, WI, GA, PA) is the confirmed
empirical claim of this paper.

\subsection{Compactness Variation Across Percentiles}

While partisan outcomes are uniform, compactness does vary across percentiles.
Table~\ref{tab:compactness} shows the normalised edge-cut fraction and mean
Polsby-Popper score for each percentile.

\begin{table}[htbp]
\centering
\caption{Normalised edge-cut fraction ($\ECnorm$) and mean Polsby-Popper
compactness by state and percentile. Lower $\ECnorm$ indicates higher
compactness. The spread $\Delta\ECnorm$ is the range from $p=0.0$ to $p=1.0$.}
\label{tab:compactness}
\begin{tabular}{lccccccr}
\toprule
\textbf{State} & \multicolumn{5}{c}{$\ECnorm$ at percentile $p$} & \textbf{Mean PP} & $\Delta\ECnorm$ \\
\cmidrule(lr){2-6}
 & $p{=}0.0$ & $p{=}0.25$ & $p{=}0.5$ & $p{=}0.75$ & $p{=}1.0$ & (at $p{=}0.0$) & \\
\midrule
NC & 0.0973 & 0.0981 & 0.0989 & 0.0997 & 0.1012 & 0.28 & 0.0039 \\
WI & 0.1142 & 0.1158 & 0.1174 & 0.1191 & 0.1208 & 0.31 & 0.0066 \\
GA & 0.1087 & 0.1104 & 0.1121 & 0.1138 & 0.1157 & 0.29 & 0.0070 \\
PA & 0.0921 & 0.0934 & 0.0948 & 0.0962 & 0.0979 & 0.32 & 0.0058 \\
TX & 0.0847 & 0.0861 & 0.0875 & 0.0890 & 0.0908 & 0.27 & 0.0061 \\
CA & 0.0712 & 0.0724 & 0.0737 & 0.0750 & 0.0765 & 0.34 & 0.0053 \\
\bottomrule
\end{tabular}
\end{table}

The compactness spread $\Delta\ECnorm$ across the full percentile range is
approximately 0.004--0.007 for all states.
This is small relative to the gap between the algorithmic plan and the
\recom{} ensemble mean (which ranges from near-zero for NC to approximately
0.030 for WI and GA).
The 101-plan bisection distribution is tightly clustered around the minimum
edge cut, consistent with B.7's finding that the coefficient of variation
of $\ECnorm$ is below 2\% for most states.

\subsection{North Carolina: The Geographic Convergence Case}
\label{subsec:nc}

North Carolina ($k = 14$) is the state where geographic convergence produces
a qualitatively different result.
As documented in G.1~\citep{g1paper}, NC's geographic structure constrains
the feasible space so tightly that the \ar{} minimum-cut plan and the \recom{}
ensemble median are nearly identical ($\Delta < 0.3\%$).

For \ps{}, this means that all five percentiles produce nearly the same plan.
The spread $\Delta\ECnorm = 0.0039$ for NC is the smallest of the six states,
confirming that geographic convergence compresses the bisection distribution.
NC's result is 7D/7R at $p = 0.0$ (the B.11 headline~\citep{b11paper}),
and this result is stable across all percentiles.

This stability has a special legal significance.
For NC, the argument ``we chose the compactness extremum'' and the argument
``we chose a plan typical of what any neutral party would draw'' lead to the
same plan.
Both compactness doctrine and representativeness doctrine converge on the
same outcome.

\subsection{Partisan Insensitivity: Seat Variation Summary}
\label{subsec:bound}

For the four states with full $T = 101$ \ps{} sweeps (NC, WI, GA, PA),
the maximum variation in projected D seats across all five percentile values
is \textbf{zero seats}.
The result is unambiguous: compactness percentile choice has no effect on
partisan outcomes for these states.

TX and CA results are extrapolated from single-run B.11 data;
full $T = 101$ sweeps for these states are reserved for the companion H.2
benchmark study.
For these two states the insensitivity claim is consistent with the B.11
point estimates but is not confirmed by actual sweep data.
The at-most-0.5-seat figure cited in earlier drafts applies conservatively
to the interpolated TX and CA entries only.

This bound is stronger than B.17's 0.3-seat bound for parameter variation,
because the bisection distribution is itself tightly concentrated (B.7 CV
$< 2\%$).
Sorting and selecting from a concentrated distribution preserves the
concentrated partisan outcome.

The result directly extends B.17's central finding: not only are partisan
outcomes insensitive to algorithm parameters, they are also insensitive to
the selection percentile within the bisection family.
The geographic structure that determines partisan outcomes (B.0, B.7) is
robust to both of these degrees of freedom.
